% Worked authority scoring example for CodeBrain paper
% Include in main.tex with: % Worked authority scoring example for CodeBrain paper
% Include in main.tex with: % Worked authority scoring example for CodeBrain paper
% Include in main.tex with: % Worked authority scoring example for CodeBrain paper
% Include in main.tex with: \input{figures/authority-example}
%
% This figure shows a small dependency graph with authority scores,
% illustrating how ROOT, DERIVED, and LEAF files differ.

\begin{figure}[t]
\centering
\begin{tikzpicture}[
    node distance=1.5cm and 2cm,
    root/.style={circle, draw, thick, fill=red!15, minimum size=1.2cm, align=center, font=\scriptsize},
    derived/.style={circle, draw, fill=yellow!15, minimum size=1.2cm, align=center, font=\scriptsize},
    leaf/.style={circle, draw, fill=green!15, minimum size=1.2cm, align=center, font=\scriptsize},
    arrow/.style={-{Stealth[length=2mm]}, gray},
    scorelabel/.style={font=\scriptsize\bfseries, text=black}
]

% ROOT nodes
\node[root] (schema) {db/\\schema};
\node[root, right=2.5cm of schema] (config) {core/\\config};

% DERIVED nodes
\node[derived, below left=1.2cm and 0.5cm of schema] (service) {services/\\user};
\node[derived, below right=1.2cm and 0.5cm of schema] (api) {api/\\users};
\node[derived, below=1.2cm of config] (auth) {services/\\auth};

% LEAF nodes
\node[leaf, below=1.2cm of service] (test1) {tests/\\user.test};
\node[leaf, below=1.2cm of api] (test2) {tests/\\api.test};

% Dependency edges (arrow = "imports from" / "depends on")
\draw[arrow] (service) -- (schema);
\draw[arrow] (api) -- (schema);
\draw[arrow] (auth) -- (schema);
\draw[arrow] (api) -- (service);
\draw[arrow] (auth) -- (config);
\draw[arrow] (service) -- (config);
\draw[arrow] (test1) -- (service);
\draw[arrow] (test2) -- (api);

% Authority scores
\node[scorelabel, above=0.05cm of schema] {$\mathcal{A}=0.92$};
\node[scorelabel, above=0.05cm of config] {$\mathcal{A}=0.87$};
\node[scorelabel, left=0.05cm of service, anchor=east] {$\mathcal{A}=0.51$};
\node[scorelabel, right=0.05cm of api, anchor=west] {$\mathcal{A}=0.42$};
\node[scorelabel, right=0.05cm of auth, anchor=west] {$\mathcal{A}=0.38$};
\node[scorelabel, below=0.05cm of test1] {$\mathcal{A}=0.05$};
\node[scorelabel, below=0.05cm of test2] {$\mathcal{A}=0.08$};

% Legend
\matrix[draw, anchor=north east, column sep=0.3cm, row sep=0.1cm, font=\scriptsize]
  at (current bounding box.south east) {
    \node[root, minimum size=0.4cm] {}; & \node{ROOT ($\geq 0.8$)}; \\
    \node[derived, minimum size=0.4cm] {}; & \node{DERIVED ($0.3$--$0.8$)}; \\
    \node[leaf, minimum size=0.4cm] {}; & \node{LEAF ($< 0.3$)}; \\
};

\end{tikzpicture}
\caption{Worked example of authority scoring on a small web application. Arrows indicate ``depends on'' relationships. The database schema and core configuration are ROOT files (high in-degree, low out-degree). Service and API files are DERIVED (moderate authority). Test files are LEAF (no dependents). A flat retrieval query for ``user management'' might rank all files containing ``user'' equally; authority-aware retrieval ranks the schema first.}
\label{fig:authority-example}
\end{figure}

%
% This figure shows a small dependency graph with authority scores,
% illustrating how ROOT, DERIVED, and LEAF files differ.

\begin{figure}[t]
\centering
\begin{tikzpicture}[
    node distance=1.5cm and 2cm,
    root/.style={circle, draw, thick, fill=red!15, minimum size=1.2cm, align=center, font=\scriptsize},
    derived/.style={circle, draw, fill=yellow!15, minimum size=1.2cm, align=center, font=\scriptsize},
    leaf/.style={circle, draw, fill=green!15, minimum size=1.2cm, align=center, font=\scriptsize},
    arrow/.style={-{Stealth[length=2mm]}, gray},
    scorelabel/.style={font=\scriptsize\bfseries, text=black}
]

% ROOT nodes
\node[root] (schema) {db/\\schema};
\node[root, right=2.5cm of schema] (config) {core/\\config};

% DERIVED nodes
\node[derived, below left=1.2cm and 0.5cm of schema] (service) {services/\\user};
\node[derived, below right=1.2cm and 0.5cm of schema] (api) {api/\\users};
\node[derived, below=1.2cm of config] (auth) {services/\\auth};

% LEAF nodes
\node[leaf, below=1.2cm of service] (test1) {tests/\\user.test};
\node[leaf, below=1.2cm of api] (test2) {tests/\\api.test};

% Dependency edges (arrow = "imports from" / "depends on")
\draw[arrow] (service) -- (schema);
\draw[arrow] (api) -- (schema);
\draw[arrow] (auth) -- (schema);
\draw[arrow] (api) -- (service);
\draw[arrow] (auth) -- (config);
\draw[arrow] (service) -- (config);
\draw[arrow] (test1) -- (service);
\draw[arrow] (test2) -- (api);

% Authority scores
\node[scorelabel, above=0.05cm of schema] {$\mathcal{A}=0.92$};
\node[scorelabel, above=0.05cm of config] {$\mathcal{A}=0.87$};
\node[scorelabel, left=0.05cm of service, anchor=east] {$\mathcal{A}=0.51$};
\node[scorelabel, right=0.05cm of api, anchor=west] {$\mathcal{A}=0.42$};
\node[scorelabel, right=0.05cm of auth, anchor=west] {$\mathcal{A}=0.38$};
\node[scorelabel, below=0.05cm of test1] {$\mathcal{A}=0.05$};
\node[scorelabel, below=0.05cm of test2] {$\mathcal{A}=0.08$};

% Legend
\matrix[draw, anchor=north east, column sep=0.3cm, row sep=0.1cm, font=\scriptsize]
  at (current bounding box.south east) {
    \node[root, minimum size=0.4cm] {}; & \node{ROOT ($\geq 0.8$)}; \\
    \node[derived, minimum size=0.4cm] {}; & \node{DERIVED ($0.3$--$0.8$)}; \\
    \node[leaf, minimum size=0.4cm] {}; & \node{LEAF ($< 0.3$)}; \\
};

\end{tikzpicture}
\caption{Worked example of authority scoring on a small web application. Arrows indicate ``depends on'' relationships. The database schema and core configuration are ROOT files (high in-degree, low out-degree). Service and API files are DERIVED (moderate authority). Test files are LEAF (no dependents). A flat retrieval query for ``user management'' might rank all files containing ``user'' equally; authority-aware retrieval ranks the schema first.}
\label{fig:authority-example}
\end{figure}

%
% This figure shows a small dependency graph with authority scores,
% illustrating how ROOT, DERIVED, and LEAF files differ.

\begin{figure}[t]
\centering
\begin{tikzpicture}[
    node distance=1.5cm and 2cm,
    root/.style={circle, draw, thick, fill=red!15, minimum size=1.2cm, align=center, font=\scriptsize},
    derived/.style={circle, draw, fill=yellow!15, minimum size=1.2cm, align=center, font=\scriptsize},
    leaf/.style={circle, draw, fill=green!15, minimum size=1.2cm, align=center, font=\scriptsize},
    arrow/.style={-{Stealth[length=2mm]}, gray},
    scorelabel/.style={font=\scriptsize\bfseries, text=black}
]

% ROOT nodes
\node[root] (schema) {db/\\schema};
\node[root, right=2.5cm of schema] (config) {core/\\config};

% DERIVED nodes
\node[derived, below left=1.2cm and 0.5cm of schema] (service) {services/\\user};
\node[derived, below right=1.2cm and 0.5cm of schema] (api) {api/\\users};
\node[derived, below=1.2cm of config] (auth) {services/\\auth};

% LEAF nodes
\node[leaf, below=1.2cm of service] (test1) {tests/\\user.test};
\node[leaf, below=1.2cm of api] (test2) {tests/\\api.test};

% Dependency edges (arrow = "imports from" / "depends on")
\draw[arrow] (service) -- (schema);
\draw[arrow] (api) -- (schema);
\draw[arrow] (auth) -- (schema);
\draw[arrow] (api) -- (service);
\draw[arrow] (auth) -- (config);
\draw[arrow] (service) -- (config);
\draw[arrow] (test1) -- (service);
\draw[arrow] (test2) -- (api);

% Authority scores
\node[scorelabel, above=0.05cm of schema] {$\mathcal{A}=0.92$};
\node[scorelabel, above=0.05cm of config] {$\mathcal{A}=0.87$};
\node[scorelabel, left=0.05cm of service, anchor=east] {$\mathcal{A}=0.51$};
\node[scorelabel, right=0.05cm of api, anchor=west] {$\mathcal{A}=0.42$};
\node[scorelabel, right=0.05cm of auth, anchor=west] {$\mathcal{A}=0.38$};
\node[scorelabel, below=0.05cm of test1] {$\mathcal{A}=0.05$};
\node[scorelabel, below=0.05cm of test2] {$\mathcal{A}=0.08$};

% Legend
\matrix[draw, anchor=north east, column sep=0.3cm, row sep=0.1cm, font=\scriptsize]
  at (current bounding box.south east) {
    \node[root, minimum size=0.4cm] {}; & \node{ROOT ($\geq 0.8$)}; \\
    \node[derived, minimum size=0.4cm] {}; & \node{DERIVED ($0.3$--$0.8$)}; \\
    \node[leaf, minimum size=0.4cm] {}; & \node{LEAF ($< 0.3$)}; \\
};

\end{tikzpicture}
\caption{Worked example of authority scoring on a small web application. Arrows indicate ``depends on'' relationships. The database schema and core configuration are ROOT files (high in-degree, low out-degree). Service and API files are DERIVED (moderate authority). Test files are LEAF (no dependents). A flat retrieval query for ``user management'' might rank all files containing ``user'' equally; authority-aware retrieval ranks the schema first.}
\label{fig:authority-example}
\end{figure}

%
% This figure shows a small dependency graph with authority scores,
% illustrating how ROOT, DERIVED, and LEAF files differ.

\begin{figure}[t]
\centering
\begin{tikzpicture}[
    node distance=1.5cm and 2cm,
    root/.style={circle, draw, thick, fill=red!15, minimum size=1.2cm, align=center, font=\scriptsize},
    derived/.style={circle, draw, fill=yellow!15, minimum size=1.2cm, align=center, font=\scriptsize},
    leaf/.style={circle, draw, fill=green!15, minimum size=1.2cm, align=center, font=\scriptsize},
    arrow/.style={-{Stealth[length=2mm]}, gray},
    scorelabel/.style={font=\scriptsize\bfseries, text=black}
]

% ROOT nodes
\node[root] (schema) {db/\\schema};
\node[root, right=2.5cm of schema] (config) {core/\\config};

% DERIVED nodes
\node[derived, below left=1.2cm and 0.5cm of schema] (service) {services/\\user};
\node[derived, below right=1.2cm and 0.5cm of schema] (api) {api/\\users};
\node[derived, below=1.2cm of config] (auth) {services/\\auth};

% LEAF nodes
\node[leaf, below=1.2cm of service] (test1) {tests/\\user.test};
\node[leaf, below=1.2cm of api] (test2) {tests/\\api.test};

% Dependency edges (arrow = "imports from" / "depends on")
\draw[arrow] (service) -- (schema);
\draw[arrow] (api) -- (schema);
\draw[arrow] (auth) -- (schema);
\draw[arrow] (api) -- (service);
\draw[arrow] (auth) -- (config);
\draw[arrow] (service) -- (config);
\draw[arrow] (test1) -- (service);
\draw[arrow] (test2) -- (api);

% Authority scores
\node[scorelabel, above=0.05cm of schema] {$\mathcal{A}=0.92$};
\node[scorelabel, above=0.05cm of config] {$\mathcal{A}=0.87$};
\node[scorelabel, left=0.05cm of service, anchor=east] {$\mathcal{A}=0.51$};
\node[scorelabel, right=0.05cm of api, anchor=west] {$\mathcal{A}=0.42$};
\node[scorelabel, right=0.05cm of auth, anchor=west] {$\mathcal{A}=0.38$};
\node[scorelabel, below=0.05cm of test1] {$\mathcal{A}=0.05$};
\node[scorelabel, below=0.05cm of test2] {$\mathcal{A}=0.08$};

% Legend
\matrix[draw, anchor=north east, column sep=0.3cm, row sep=0.1cm, font=\scriptsize]
  at (current bounding box.south east) {
    \node[root, minimum size=0.4cm] {}; & \node{ROOT ($\geq 0.8$)}; \\
    \node[derived, minimum size=0.4cm] {}; & \node{DERIVED ($0.3$--$0.8$)}; \\
    \node[leaf, minimum size=0.4cm] {}; & \node{LEAF ($< 0.3$)}; \\
};

\end{tikzpicture}
\caption{Worked example of authority scoring on a small web application. Arrows indicate ``depends on'' relationships. The database schema and core configuration are ROOT files (high in-degree, low out-degree). Service and API files are DERIVED (moderate authority). Test files are LEAF (no dependents). A flat retrieval query for ``user management'' might rank all files containing ``user'' equally; authority-aware retrieval ranks the schema first.}
\label{fig:authority-example}
\end{figure}
